\chapter{Smart-BMS Software-in-the-Loop Validation}
\label{chap:sil_validation}

\section{Coupling the Physical and Algorithmic Domains}
\label{sec:sil_coupling}

Chapter~10 developed the Simulink-based physical plant model of the battery cell -- a first-order Thevenin ECM parameterized with the offline and online identification methods of Chapters~4 through~9 -- and subjected it to a structured dynamic stress-testing framework spanning idle, partial-load, and high-current excitation. That plant model, on its own, answers only half of the question this report set out to address: it shows how a battery \emph{behaves}, but it says nothing about how accurately a BMS's parameter and state estimation algorithms -- the RLS, EKF, UKF, observer-based, co-estimation, and AI-driven methods developed across Chapters~4 through~9 -- \emph{track} that behavior once it is corrupted by measurement noise, initialized from imperfect starting conditions, and constrained by the limited computational resources of an embedded microcontroller. This chapter closes that gap by coupling the physical plant model of Chapter~10 to the algorithmic (estimation) domain in a software-in-the-loop (SIL) simulation environment, and by evaluating the resulting closed-loop system against the transient tracking and real-time computational criteria that ultimately determine whether a given estimation method, as developed mathematically in earlier chapters, is fit for deployment in a production BMS.

\subsection{Chapter Roadmap}
\label{ssec:ch11_roadmap}

Section~\ref{ssec:xil_taxonomy} places SIL testing within the broader MIL/SIL/PIL/HIL verification hierarchy and Section~\ref{ssec:sil_architecture} develops the specific signal architecture used to couple the Chapter~10 plant model to the algorithmic domain, with Sections~\ref{ssec:sil_nonrealtime}--\ref{ssec:sil_reproducibility} addressing the practical considerations -- non-real-time execution, test-case coverage, and statistical reproducibility -- that govern how the results of this chapter should be generated and interpreted. Section~\ref{sec:transient_tracking} develops the transient-tracking evaluation itself, fixing the quantitative metrics used throughout the chapter (Section~\ref{ssec:transient_metrics}) before applying them to the abrupt load-disconnect scenario (Section~\ref{ssec:load_disconnect}) and the related voltage-estimate stabilization problem (Section~\ref{ssec:voltage_peak_stabilization}). Section~\ref{sec:realtime_efficiency} shifts from accuracy to computational cost, establishing the real-time execution-rate target (Section~\ref{ssec:realtime_requirements}), comparing the algorithm families of Chapters~5 through~9 on both qualitative and reported-hardware bases (Section~\ref{ssec:algorithm_complexity}), and cataloguing implementation-level optimization techniques (Section~\ref{ssec:embedded_optimization}) together with the memory-footprint and pack-scalability considerations that extend beyond raw execution time (Sections~\ref{ssec:memory_footprint}--\ref{ssec:pack_scalability}). Section~\ref{sec:summary_ch11} closes the chapter by synthesizing the accuracy/cost trade-off this chapter has developed and bridging to the concluding synthesis of Chapter~12.

\subsection{Notation Used in This Chapter}
\label{ssec:ch11_notation}

Table~\ref{tab:ch11_notation} collects the symbols introduced specifically in this chapter, extending the notation table first established in Chapter~2, Section~2.1.4 and reused without redefinition, per the cross-chapter consistency requirement of Chapter~1, throughout Chapters~3 through~10.

\begin{table}[htbp]
	\centering
	\caption{Notation introduced in Chapter~11.}
	\label{tab:ch11_notation}
	\begin{tabular}{p{2.2cm}p{7.2cm}p{2.4cm}}
		\toprule
		Symbol & Meaning & First used \\
		\midrule
		$\hat{x}(t)$ & Estimated value of a state or parameter $x(t)$ & Eq.~\eqref{eq:rmse} \\
		$\mathrm{RMSE}$ & Root-mean-square error between a true and estimated trajectory & Eq.~\eqref{eq:rmse} \\
		$\hat{V}_t(t)$ & Estimator-internal predicted terminal voltage & Sec.~\ref{ssec:voltage_peak_stabilization} \\
		$N_{\text{trial}}$ & Number of independent SIL trials in a repeated-trial evaluation & Sec.~\ref{ssec:sil_reproducibility} \\
		$n$ & Dimension of the estimated state/parameter vector & Sec.~\ref{ssec:algorithm_complexity} \\
		\bottomrule
	\end{tabular}
\end{table}

\subsection{Verification and Validation Levels: MIL, SIL, PIL, and HIL}
\label{ssec:xil_taxonomy}

Before developing the SIL architecture used in this chapter, it is useful to place it within the broader verification and validation (V\&V) hierarchy used in automotive and embedded-systems engineering practice, since the choice of validation level determines both what kinds of defects a given test can expose and what it costs, in engineering time and hardware, to run that test. This hierarchy is commonly organized around four progressively more hardware-realistic stages, mapped onto the right-hand (verification) branch of the V-model of embedded software development \cite{patsnap_hil_sil}:

\begin{itemize}
	\item \textbf{Model-in-the-Loop (MIL).} Both the estimation algorithm and the plant model are executed as high-level Simulink models, with no compiled code and no target hardware involved; computations are performed in floating-point arithmetic and the resulting signals serve as the reference values against which later, more hardware-realistic stages are compared \cite{empirical_mil_sil_arxiv}. The Chapter~10 plant model, exercised directly in Simulink against a Simulink-native implementation of an estimation algorithm, constitutes MIL-level validation.
	\item \textbf{Software-in-the-Loop (SIL).} The estimation algorithm is compiled (or auto-code-generated, in the MathWorks Embedded Coder sense) into the same source-code form that will ultimately run on the BMS microcontroller, and this compiled code is executed on a host PC, still coupled to the plant model rather than to physical hardware \cite{mathworks_bms_solutions}. SIL testing is the primary validation methodology used throughout this chapter: it allows the estimation algorithms of Chapters~5 through~9 to be evaluated in their deployment-representative, compiled numerical form, while retaining the convenience, repeatability, and diagnostic access of a fully simulated environment.
	\item \textbf{Processor-in-the-Loop (PIL).} The same compiled code is moved onto the actual target microcontroller (or an equivalent evaluation board), executed in isolation, with its inputs and outputs still driven by and returned to the plant model running on a host PC \cite{embitel_mil_sil_hil}. PIL testing exposes target-specific numerical behavior -- fixed-point rounding, compiler optimization effects, and genuine execution timing -- that a host-PC SIL simulation cannot fully replicate, and is the natural extension of the real-time computational efficiency analysis of Section~\ref{sec:realtime_efficiency} onto physical target hardware.
	\item \textbf{Hardware-in-the-Loop (HIL).} The complete BMS hardware -- including its analog front-end, current and voltage sensing circuitry, and communication interfaces -- is connected to a real-time simulated battery pack, typically implemented on dedicated real-time target hardware capable of emulating individual cell voltages and executing the plant model at a fixed real-time step \cite{mathworks_hil_bms_webinar}. HIL represents the final pre-vehicle validation stage and is generally reported to require battery plant model execution at update rates on the order of 1~kHz to comfortably exceed the 100--500~Hz execution rate typical of production BMS estimation algorithms \cite{mathworks_hil_bms_video}.
\end{itemize}

\begin{table}[htbp]
	\centering
	\caption{Verification and validation levels for BMS estimation algorithms, in order of increasing hardware fidelity.}
	\label{tab:xil_levels}
	\begin{tabular}{p{1.6cm}p{3.0cm}p{3.4cm}p{3.4cm}}
		\toprule
		Level & Estimation algorithm form & Plant model & Primary purpose \\
		\midrule
		MIL & Simulink model (floating point) & Simulink model & Algorithmic correctness, reference generation \\
		SIL & Compiled/auto-generated code, host PC & Simulink model, host PC & This chapter: transient tracking and preliminary timing \\
		PIL & Compiled code, target microcontroller & Simulink model, host PC & Target-specific numerics and execution timing \\
		HIL & Deployed BMS hardware & Real-time target (e.g. Speedgoat) & Full closed-loop, sensor-realistic validation \\
		\bottomrule
	\end{tabular}
\end{table}

Progression through these levels is understood, in standard automotive V\&V practice, to be sequential rather than a matter of selecting a single preferred level: a defect caught at MIL costs comparatively little to fix, while the same defect surfacing for the first time at HIL, or later, can require a substantially more costly re-design \cite{askui_xil_hierarchy}. This report's scope -- parameter and state estimation algorithm development and validation, rather than full ECU hardware qualification -- places its natural validation ceiling at SIL, with PIL and HIL noted here as the logical next steps beyond this report's boundary, consistent with the scope established in Chapter~1.

\subsection{SIL Architecture for This Report}
\label{ssec:sil_architecture}

The SIL environment developed in this chapter couples two previously independent halves of this report into a single closed-loop simulation:

\begin{enumerate}
	\item \textbf{The physical domain}, consisting of the Chapter~10 Simulink plant model -- the first-order Thevenin ECM (Chapter~2, Equations~2.5--2.9) parameterized from the offline and online identification results of Chapters~4 through~9, and exercised under the dynamic stress-testing current profiles developed in Chapter~10, Section~10.3.
	\item \textbf{The algorithmic domain}, consisting of a compiled, deployment-representative implementation of one or more estimation algorithms drawn from Chapters~5 through~9 -- for example, an EKF-based joint state/parameter estimator (Chapter~5, Chapter~7) or a sliding-mode observer (Chapter~6) -- that receives the same current excitation applied to the plant model, together with a noise-corrupted version of the plant's terminal voltage output (using the sensor noise characterization of Chapter~3, Section~3.5), and produces estimates of SOC and/or the ECM parameters.
\end{enumerate}

The essential feature that distinguishes this SIL coupling from either domain in isolation is the availability, within the simulation, of \emph{ground truth}: because the plant model's internal states ($z(t)$, $U_p(t)$) and parameters ($R_0$, $R_1$, $C_1$) are directly accessible signals within the Simulink model, the estimation algorithm's output can be compared directly against the true simulated quantity at every time step, something that is not possible with a physical cell, whose true internal SOC and parameter values are never directly measurable. This ground-truth availability is what makes SIL simulation a uniquely efficient environment for quantifying the transient tracking performance developed in Section~\ref{sec:transient_tracking}, ahead of the comparatively more expensive PIL and HIL stages noted in Section~\ref{ssec:xil_taxonomy}.

\begin{quote}
\textbf{Signal flow of the SIL coupling used in this chapter:}
\begin{enumerate}
	\item A current excitation profile $I(t)$ (drawn from the dynamic stress-testing framework of Chapter~10, Section~10.3) is applied identically to both the plant model and, where a current sensor model is included, to the algorithmic domain's current input.
	\item The plant model integrates its internal state equations (Chapter~2, Equation~2.9) and produces the true terminal voltage $V_t(t)$ and the true internal states $z(t)$, $U_p(t)$.
	\item A measurement model, using the noise characterization of Chapter~3, Section~3.5, corrupts $V_t(t)$ (and, where relevant, $I(t)$) with representative sensor noise, bias, and quantization before this signal reaches the algorithmic domain.
	\item The estimation algorithm under test (Chapters~5--9) processes the noisy current and voltage measurements and produces estimates $\hat{z}(t)$, $\hat{R}_0(t)$, $\hat{R}_1(t)$, $\hat{C}_1(t)$ (as applicable to the algorithm's design).
	\item The estimation error signals -- $z(t) - \hat{z}(t)$ and the corresponding parameter errors -- are logged alongside the true and estimated terminal voltage, and are the primary output evaluated in Sections~\ref{sec:transient_tracking} and~\ref{sec:realtime_efficiency}.
\end{enumerate}
\end{quote}

\subsection{Non-Real-Time Execution and Its Implications}
\label{ssec:sil_nonrealtime}

A property of SIL simulation that is easy to overlook, but that materially affects how the results of this chapter should be interpreted, is that SIL execution is not, in general, constrained to run at the same rate as the physical process it simulates: unlike HIL, where the plant model must execute on dedicated real-time hardware fast enough to keep pace with wall-clock time and with the connected physical BMS hardware, a SIL simulation can run faster or slower than real time, limited only by host-PC computational throughput \cite{opalrt_sil}. This is directly advantageous for the transient tracking evaluation of Section~\ref{sec:transient_tracking}, where long-duration dynamic stress profiles can be simulated and their estimation error statistics computed without waiting for wall-clock-real-time execution. It is, however, precisely the property that must be set aside when the real-time computational efficiency analysis of Section~\ref{sec:realtime_efficiency} is performed: because SIL execution time on a host PC bears no fixed relationship to execution time on the target microcontroller, per-step timing measurements taken in SIL are indicative only in a coarse, relative sense, and the report's more specific timing figures in Section~\ref{sec:realtime_efficiency} are drawn from published PIL/HIL and direct embedded-hardware implementation studies rather than from host-PC SIL timing alone.

\subsection{Test Case Design and Coverage}
\label{ssec:sil_test_design}

A single load-disconnect transient, however carefully instrumented, is not sufficient on its own to characterize an estimation algorithm's behavior across the operating envelope a production BMS must handle. Consistent with the dynamic stress-testing framework already established in Chapter~10, Section~10.3, the SIL test suite used in this chapter is organized around a small number of current-excitation archetypes, each targeting a distinct aspect of estimator behavior:

\begin{itemize}
	\item \textbf{Idle and partial-load transients} (Chapter~10, Section~10.3.1), which exercise the estimator's behavior under small, low-rate current excursions and are primarily used to verify steady-state accuracy (via the RMSE metric of Section~\ref{ssec:transient_metrics}) rather than transient tracking specifically.
	\item \textbf{Abrupt load disconnects} (Section~\ref{ssec:load_disconnect}), which exercise the fastest dynamics retained in the ECM structures of Chapter~2 and are the primary scenario used to expose the three failure modes catalogued in Table~\ref{tab:load_disconnect_modes}.
	\item \textbf{High-current stress and thermal boundary excitation} (Chapter~10, Section~10.3.2), which, when combined with the thermal-coupled model extensions of Chapter~2, Section~2.4, additionally exercises the estimator's robustness to the temperature-dependent parameter drift discussed there, and is used to verify that an estimator tuned under nominal-temperature conditions does not diverge when the plant's true parameters shift with self-heating.
	\item \textbf{Standardized dynamic drive-cycle profiles}, of the kind already used for offline ECM validation in Chapter~2, Section~2.2.5 and Chapter~3, Section~3.3.5, repurposed here as a SIL-level closed-loop test rather than an open-loop model-fit validation, so that estimator tracking accuracy is assessed under application-representative excitation rather than only under the idealized step and pulse inputs of the preceding scenarios.
\end{itemize}

This layered test design -- from simple, diagnostic step transients through to full drive-cycle excitation -- mirrors standard test-coverage practice in embedded software V\&V, where structured, individually diagnostic test cases are used to isolate specific failure modes before broader, application-representative scenarios are used to confirm overall fitness for purpose \cite{empirical_mil_sil_arxiv}.

\subsection{Reproducibility and Statistical Evaluation}
\label{ssec:sil_reproducibility}

Because the measurement noise injected into the SIL loop (Section~\ref{ssec:sil_architecture}) is, by construction, a random process characterized statistically in Chapter~3, Section~3.5.2, a single SIL run under a given current excitation profile represents only one realization of that noise process, and the resulting RMSE, convergence-time, and overshoot metrics (Section~\ref{ssec:transient_metrics}) are themselves random variables rather than fixed quantities. Sound SIL evaluation practice therefore repeats each test scenario across multiple independent noise realizations -- a Monte Carlo approach directly analogous to the repeated-trial evaluation practice common in the broader Kalman-filtering and observer-design literature referenced in Chapters~5 and~6 -- and reports the resulting metric \emph{distributions} (for instance, mean and standard deviation of RMSE across trials) rather than a single-run point value. This distinction matters in particular for the overshoot and convergence-time metrics of Section~\ref{ssec:transient_metrics}, whose single-run values can be disproportionately influenced by whether a particular noise realization happens to align constructively or destructively with the transient event under test; a stabilization technique that reduces \emph{mean} overshoot across repeated trials, evaluated as described in Section~\ref{ssec:voltage_peak_stabilization}, provides materially stronger evidence than a single favorable-looking run.

\section{Transient Tracking Performance}
\label{sec:transient_tracking}

Validation of a model's structural adequacy -- the question of whether a given ECM order or electrochemical model captures enough physical detail to reproduce a cell's measured behavior -- was addressed, for the model families of Chapter~2, using the offline validation principle of Chapter~2, Section~2.2.5: simulate the identified model against a current profile not used in parameter extraction, and compare against independently measured voltage. That principle answers whether the \emph{model} is adequate. It does not, on its own, answer whether an \emph{estimation algorithm} coupled to that model, and constrained to operate on noisy, real-time measurements rather than clean offline characterization data, can track the model's internal states and parameters with acceptable transient accuracy. This section develops that second, estimator-focused validation question, using the SIL coupling architecture of Section~\ref{sec:sil_coupling} to make the plant's ground-truth trajectories available for direct comparison against the estimator's output -- an evaluation that, as noted in Section~\ref{ssec:sil_architecture}, is uniquely tractable at the SIL validation level.

\subsection{Evaluation Metrics}
\label{ssec:transient_metrics}

Before presenting the two specific transient scenarios developed in Sections~\ref{ssec:load_disconnect} and~\ref{ssec:voltage_peak_stabilization}, it is necessary to fix the quantitative metrics by which "tracking performance" is judged throughout this chapter, since the qualitative discussion of estimator behavior that follows is grounded in these metrics rather than in subjective visual comparison of plotted trajectories.

\begin{itemize}
	\item \textbf{Root-mean-square error (RMSE).} For a state or parameter estimate $\hat{x}(t)$ compared against the plant's ground-truth value $x(t)$ over a simulation window of $N$ samples,
	\begin{equation}
		\mathrm{RMSE} = \sqrt{\frac{1}{N}\sum_{k=1}^{N} \big(x_k - \hat{x}_k\big)^2},
		\label{eq:rmse}
	\end{equation}
	is the standard aggregate accuracy metric used throughout the estimation literature referenced in Chapters~5 through~7, and is used here as the primary steady-state (post-transient) accuracy indicator.
	\item \textbf{Convergence time.} The elapsed time, following a transient excitation event (a current step, or filter initialization from an incorrect initial condition), required for the estimation error to fall within and remain within a specified tolerance band -- commonly expressed as a percentage of the corresponding steady-state value, analogous to the settling-time definition used in classical control system step-response analysis.
	\item \textbf{Overshoot.} The maximum excursion of an estimate beyond its eventual steady-state value following a transient, expressed either in absolute terms or as a percentage of the transient's magnitude; overshoot in the \emph{estimated} voltage or state is distinct from, but can be induced by, filter gain or observer gain tuning, as developed in Section~\ref{ssec:voltage_peak_stabilization}.
	\item \textbf{Estimation error during the transient window itself}, as distinct from post-transient RMSE, since a filter or observer that converges to an accurate steady-state value can nonetheless exhibit a large, transient-specific error spike immediately following a step change in current -- the specific failure mode that Sections~\ref{ssec:load_disconnect} and~\ref{ssec:voltage_peak_stabilization} are designed to expose.
\end{itemize}

\subsection{Response to Abrupt Load Disconnects}
\label{ssec:load_disconnect}

\subsubsection{Test Scenario}

An abrupt load disconnect -- a step change in applied current from a sustained nonzero value (discharge or charge) to zero, applied without ramping -- is among the most demanding transient excitations an estimation algorithm can be asked to track, precisely because it exercises the fastest dynamics retained in the ECM structures of Chapter~2 without any of the smoothing that a ramped or filtered current profile would provide. This scenario is a direct SIL-domain counterpart to the idle and partial-load stress-testing scenarios developed in the physical plant model of Chapter~10, Section~10.3.1, now evaluated from the perspective of the estimation algorithm coupled to that plant rather than the plant's open-loop response alone.

\subsubsection{Expected Plant Response}

The plant model's true terminal voltage response to a load-disconnect step follows directly from the Thevenin step-response relation already derived in Chapter~2, Equation~2.13: at the instant of disconnect, the terminal voltage jumps discontinuously by $I_0 R_0$ (the ohmic term vanishing instantaneously along with the current), after which the polarization voltage $U_p(t)$ relaxes exponentially toward zero with time constant $\tau_1 = R_1 C_1$, so that the terminal voltage approaches its new open-circuit value $V_{oc}(z)$ asymptotically rather than instantaneously. This two-part response -- an instantaneous step followed by a first-order exponential relaxation -- is precisely the signature that the offline HPPC extraction procedure of Chapter~3, Section~3.3 is designed to characterize, and it recurs here as the reference trajectory against which the \emph{estimator's} response, not merely the plant's own response, is now evaluated.

\subsubsection{Estimator Tracking Behavior}

The central question this scenario is designed to answer is not whether the plant model reproduces the expected step-plus-relaxation voltage trajectory -- that behavior is guaranteed by construction, given the ECM structure of Chapter~2 -- but whether the estimation algorithm coupled to it, receiving only the noisy current and voltage measurements, can track the corresponding \emph{state} and \emph{parameter} trajectories with acceptably small transient error. Three distinct failure modes are of particular concern in this scenario, each traceable to specific mathematical properties of the estimation methods developed in earlier chapters:

\begin{enumerate}
	\item \textbf{Delayed convergence following the step.} A Kalman-filter-based estimator (Chapter~5) with poorly tuned process-noise covariance $Q$ relative to the true rate of state variation can under-respond to a step change, producing a state estimate that lags the true trajectory during the fast initial transient before eventually converging; this is a direct manifestation of the persistent-excitation and observability considerations flagged qualitatively in Chapter~2, Section~2.2.3 and developed formally as an identifiability concern in Chapter~4.
	\item \textbf{Transient parameter estimate corruption.} For a joint or co-estimation architecture (Chapter~7) in which SOC and ECM parameters are estimated simultaneously, an abrupt current step -- being a strong, sudden excitation -- can produce a transient in the \emph{parameter} estimates (particularly $R_0$, whose associated signal is directly proportional to the step itself) even though the true parameter values have not changed; this is a known consequence of the weak, and momentarily correlated, distinguishability between state and parameter sensitivities during a sharp transient, a concern raised in general terms in Chapter~7, Section~7.1.
	\item \textbf{Observer or filter divergence under model mismatch.} Where the ECM order used by the estimator (e.g. a first-order Thevenin structure) does not fully capture the plant's true dynamics (for instance, if the plant model itself retains second-order polarization behavior not represented in the estimator's internal model), the abrupt load-disconnect transient is the excitation most likely to expose this mismatch, since it drives the largest and fastest-changing overpotential of any scenario in the dynamic stress-testing framework of Chapter~10.
\end{enumerate}

\begin{table}[htbp]
	\centering
	\caption{Transient behaviors evaluated under the abrupt load-disconnect scenario, and their traceable cause within the estimator architectures of earlier chapters.}
	\label{tab:load_disconnect_modes}
	\begin{tabular}{p{3.4cm}p{4.6cm}p{4.6cm}}
		\toprule
		Observed behavior & Likely cause & Related report section \\
		\midrule
		Lagging state estimate & Process-noise covariance mistuned relative to true transient rate & Ch.~5 (Kalman filter tuning) \\
		Transient parameter estimate spike & Weak state/parameter distinguishability during sharp excitation & Ch.~7, Sec.~7.1 (coupling of states and parameters) \\
		Growing or oscillatory error & ECM order mismatch between estimator and plant, or unstable observer gain & Ch.~2 (model order); Ch.~6 (robustness and convergence) \\
		\bottomrule
	\end{tabular}
\end{table}

\subsubsection{Interpreting SIL Results for This Scenario}

Because the plant's true state and parameter trajectories are directly available in SIL simulation (Section~\ref{ssec:sil_architecture}), the load-disconnect scenario is particularly well suited to distinguishing, empirically, between the three failure modes above: a lagging-but-eventually-converging state estimate, a parameter estimate that spikes transiently but returns to its correct value, and a genuinely diverging or persistently biased estimate are each visually and statistically distinct in the logged error signals of Section~\ref{ssec:transient_metrics}, even though all three might otherwise be reported, in a less carefully instrumented test, simply as "estimation error." This diagnostic value is one of the specific reasons SIL testing is positioned, in the V\&V hierarchy of Section~\ref{ssec:xil_taxonomy}, ahead of PIL and HIL: the additional visibility into internal plant states that SIL provides is progressively lost at later validation stages, where only externally measurable signals remain accessible.

\subsection{Dynamic Stabilization of Voltage Peaks}
\label{ssec:voltage_peak_stabilization}

\subsubsection{Motivation}

Where Section~\ref{ssec:load_disconnect} addressed the \emph{estimator's} response to a transient in the \emph{true} system, this section addresses a related but distinct concern: the estimator's own predicted terminal voltage -- the quantity $\hat{V}_t(t)$ computed internally by a Kalman filter or observer as part of its innovation/residual calculation (Chapters~5 and~6) -- can itself exhibit transient peaks or oscillatory behavior that do not correspond to any physical voltage excursion of the plant, but are instead artifacts of filter or observer gain tuning. Left unaddressed, such artifacts can propagate into downstream BMS functions that consume the estimator's output directly, including the state-of-power calculations that build on the HPPC-derived power-capability characterization of Chapter~3, Section~3.3.4.

\subsubsection{Sources of Voltage-Estimate Oscillation}

Two mechanisms, both traceable to material already developed in Chapters~5 and~6, are the principal sources of this behavior:

\begin{enumerate}
	\item \textbf{Excessive Kalman gain following a large innovation.} Immediately after a load-disconnect-type transient (Section~\ref{ssec:load_disconnect}), the measurement residual (innovation) used by an EKF or UKF (Chapter~5) is, by construction, large; if the filter's gain is not appropriately bounded or adaptively reduced, this large innovation can be over-weighted in the subsequent state update, producing a correction that overshoots the true state before settling back -- an estimator-domain analogue of classical control overshoot, governed by the same underlying trade-off between convergence speed and stability margin that recurs throughout the observer robustness and convergence analysis of Chapter~6, Section~6.4.
	\item \textbf{Observer gain selection in PI and sliding-mode observers.} For the proportional-integral (PI) and sliding-mode observer (SMO) architectures of Chapter~6, the proportional/switching gain directly trades off convergence speed against susceptibility to chattering or overshoot in the presence of measurement noise (Chapter~3, Section~3.5): a high gain, chosen to minimize convergence time under the load-disconnect scenario of Section~\ref{ssec:load_disconnect}, correspondingly increases the observer's sensitivity to the sensor noise characterized in Chapter~3, potentially reintroducing the very voltage-estimate oscillation this section is concerned with.
\end{enumerate}

\subsubsection{Mathematical Origin of Gain-Induced Overshoot}

The overshoot mechanism described qualitatively above can be stated more precisely using the standard Kalman filter measurement-update equation already introduced structurally in Chapter~5. Given a prior state estimate $\hat{x}_k^-$ and measurement $y_k$, the updated estimate is

\begin{equation}
	\hat{x}_k = \hat{x}_k^- + K_k \big(y_k - h(\hat{x}_k^-)\big),
	\label{eq:kf_update}
\end{equation}

\noindent where $K_k$ is the Kalman gain and the bracketed term is the innovation -- the mismatch between the actual measurement and the value predicted from the prior state estimate via the measurement function $h(\cdot)$ (for an ECM-based estimator, $h(\cdot)$ is the terminal-voltage relation of Chapter~2, Equation~2.7 or~2.11). Immediately following a load-disconnect transient (Section~\ref{ssec:load_disconnect}), the innovation is large simply because the prior estimate, propagated from before the step, has not yet accounted for the new operating point; if $K_k$ is large at this moment -- as it will tend to be immediately after a filter's process-noise covariance $Q$ has been (correctly) sized to allow fast tracking of genuine transients -- the correction term $K_k(y_k - h(\hat{x}_k^-))$ can overcorrect, driving $\hat{x}_k$ past the true post-transient value before the filter's subsequent cycles pull it back. This is the same fast-tracking-versus-noise-sensitivity trade-off already discussed, at the level of $Q$ and $R$ matrix tuning, in Chapter~5, and it is restated here specifically in terms of its observable consequence in the SIL-logged voltage-estimate trajectory: a large, correctly sized $Q$ that enables fast convergence in the sense of Section~\ref{ssec:transient_metrics} is, mechanistically, the same design choice that produces the overshoot this section addresses, so that the two cannot be optimized independently of one another.

\subsubsection{Stabilization Approaches}

Consistent with the robustness and convergence analysis already developed in Chapter~6, Section~6.4, and the adaptive process-noise tuning discussion of Chapter~5, three general stabilization strategies are relevant to suppressing spurious voltage-estimate peaks in the SIL environment of this chapter, without repeating the full algorithmic derivations already given in those chapters:

\begin{itemize}
	\item \textbf{Adaptive innovation-based gain scaling}, in which the filter's effective gain is reduced when the innovation magnitude exceeds a statistically expected bound (given the measurement-noise characterization of Chapter~3, Section~3.5.2), directly limiting the overshoot mechanism identified above without requiring a full re-derivation of the filter's covariance propagation.
	\item \textbf{Gain scheduling or boundary-layer smoothing for sliding-mode observers}, replacing a hard switching term with a smoothed (saturation-function) approximation within a boundary layer around the sliding surface, a standard technique for reducing chattering-induced oscillation while retaining the fast convergence property that motivates the sliding-mode approach in the first place (Chapter~6, Section~6.3).
	\item \textbf{Post-estimation low-pass filtering of the voltage-estimate residual}, applied cautiously and only to the diagnostic/monitoring output rather than to the feedback path of the estimator itself, so as not to reintroduce the fast-transient suppression risk already flagged for raw sensor data in Chapter~3, Section~3.5.3.
\end{itemize}

\subsubsection{SIL-Based Evaluation of Stabilization Effectiveness}

Because the SIL environment of this chapter provides direct access to both the plant's true voltage trajectory and the estimator's internal predicted voltage (Section~\ref{ssec:sil_architecture}), the effectiveness of any of the stabilization approaches above can be quantified directly as a reduction in the overshoot and convergence-time metrics defined in Section~\ref{ssec:transient_metrics}, evaluated specifically over the window immediately following a load-disconnect event. This before/after comparison -- estimator behavior with and without a given stabilization measure applied, under an identical current excitation and identical plant ground truth -- is a validation exercise that is, in practice, difficult to perform cleanly outside a SIL (or MIL) environment, since it requires the kind of repeatable, ground-truth-referenced comparison that physical hardware testing cannot directly provide.

\subsubsection{Illustrative SIL Trial (Invented Numerical Values, for Demonstration Only)}
\label{sssec:sil_illustrative_example}

To make the metrics of Section~\ref{ssec:transient_metrics} and the failure modes of Table~\ref{tab:load_disconnect_modes} concrete, it is useful to walk through a single hypothetical SIL trial illustrating how the reported quantities would be read and interpreted, using invented numerical values consistent with the illustrative-example convention already established in Chapter~2, Sections~2.2.2 and~2.4.2. Suppose a SIL run applies a load-disconnect step from $I_0 = 20\ \text{A}$ discharge to $0\ \text{A}$ at $t = 0$, to a plant model parameterized with $R_0 = 12\ \text{m}\Omega$, $R_1 = 8\ \text{m}\Omega$, $\tau_1 = 15\ \text{s}$ (Chapter~2, Section~2.2.2), against an EKF-based estimator (Chapter~5) built on the same first-order Thevenin structure. A hypothetical post-processing of the logged trial might report an SOC-estimate convergence time of $4.2\ \text{s}$ to within a 2\% tolerance band, a peak transient SOC error of $0.6\%$ occurring at $t \approx 1.1\ \text{s}$, and a steady-state (post-transient) SOC RMSE of $0.15\%$ over the subsequent 60-second window, together with a transient $\hat{R}_0$ estimate spike from its converged value of $12\ \text{m}\Omega$ up to a hypothetical $14.3\ \text{m}\Omega$ before relaxing back within approximately $3\ \text{s}$ -- an example of the transient parameter-estimate corruption mode identified in Table~\ref{tab:load_disconnect_modes}. As with every invented numerical example in this report, these figures are illustrative only, intended solely to demonstrate how the metrics of Section~\ref{ssec:transient_metrics} would be read from an actual SIL log, and must not be cited as representative performance figures for any specific algorithm, cell, or implementation.

\subsection{Statistical Evaluation Across Multiple Trials}
\label{ssec:sil_monte_carlo}

Following the reproducibility principle established in Section~\ref{ssec:sil_reproducibility}, the illustrative single-trial result of Section~\ref{sssec:sil_illustrative_example} would, in a properly conducted evaluation, be repeated across a batch of independent noise realizations, with the resulting RMSE, convergence-time, and overshoot values summarized as a distribution (for example, via their sample mean and standard deviation, or a boxplot across trials) rather than reported as a single number. Table~\ref{tab:sil_metric_summary_template} shows the structure of such a summary table, populated with placeholder entries rather than invented numerical values, since no specific algorithm implementation or trial batch has been executed as part of this report; the table is included to fix the reporting format expected of any SIL evaluation carried out using the framework of this chapter, consistent with the validation-methodology emphasis stated in Chapter~1, Section~1.3.

\begin{table}[htbp]
	\centering
	\caption{Template for reporting SIL transient-tracking results across repeated trials (values to be populated from an actual test campaign; not populated here).}
	\label{tab:sil_metric_summary_template}
	\begin{tabular}{p{3.4cm}p{2.6cm}p{2.6cm}p{2.6cm}p{2.4cm}}
		\toprule
		Scenario & Metric & Mean & Std.\ dev. & $N$ trials \\
		\midrule
		Load disconnect & SOC convergence time (s) & -- & -- & -- \\
		Load disconnect & Peak transient SOC error (\%) & -- & -- & -- \\
		Load disconnect & Post-transient SOC RMSE (\%) & -- & -- & -- \\
		Load disconnect & $\hat{R}_0$ transient overshoot (m$\Omega$) & -- & -- & -- \\
		High-current stress & SOC RMSE (\%) & -- & -- & -- \\
		Drive-cycle profile & SOC RMSE (\%) & -- & -- & -- \\
		\bottomrule
	\end{tabular}
\end{table}

\subsection{Fault Injection Testing}
\label{ssec:fault_injection}

Beyond the nominal-operation transient scenarios of Section~\ref{ssec:sil_test_design}, the SIL environment developed in this chapter is also well suited to fault-injection testing, in which the noise and bias models of Chapter~3, Section~3.5 are deliberately extended beyond their nominal statistical characterization to include the sensor fault modes previewed there -- stuck values, abrupt offset shifts, and signal dropout. Because the plant's true state remains available as ground truth throughout a SIL simulation (Section~\ref{ssec:sil_architecture}), fault-injection SIL testing can directly quantify not only whether a given estimation algorithm's accuracy degrades gracefully under a sensor fault, but also whether the model-based residual monitoring approach previewed in Chapter~3, Section~3.5.5 correctly flags the fault within an acceptable detection latency, without producing an excessive rate of false positives under ordinary sensor noise. This form of testing is a natural complement to the transient-tracking evaluation of Section~\ref{sec:transient_tracking} -- both exercise the estimator under a deliberately adverse but well-characterized input -- and it extends the accuracy/cost synthesis of this chapter with a third axis, robustness to sensor-level faults, that is particularly relevant to the safety-motivated parameter-estimation rationale first established in Chapter~1, Section~1.3.1.

\section{Real-Time Computational Efficiency}
\label{sec:realtime_efficiency}

Sections~\ref{sec:sil_coupling} and~\ref{sec:transient_tracking} evaluated estimation algorithms exclusively on accuracy grounds -- how closely a given method's output tracks the plant's true states and parameters under transient excitation. Accuracy alone, however, does not determine whether a given method can be deployed in a production BMS: every estimation algorithm developed in Chapters~5 through~9 must ultimately execute within the fixed, and often severely constrained, computational budget of an automotive- or consumer-grade microcontroller, repeated at every sample interval for the operational lifetime of the pack. This section develops the computational side of that deployment question, first establishing the execution-rate target against which cost is judged (Section~\ref{ssec:realtime_requirements}), then comparing the algorithm families of earlier chapters on both structural and reported-hardware grounds (Section~\ref{ssec:algorithm_complexity}), and finally cataloguing the implementation-level techniques -- and the memory and pack-scale considerations that accompany them -- by which a selected algorithm's computational footprint can be reduced without altering its underlying mathematics (Sections~\ref{ssec:embedded_optimization}--\ref{ssec:pack_scalability}).

\subsection{Execution Rate Requirements}
\label{ssec:realtime_requirements}

Before comparing the computational cost of individual estimation algorithms, it is useful to state the execution-rate target against which that cost is judged. Production BMS estimation algorithms are commonly reported to execute at update rates in the range of 100~Hz to 500~Hz \cite{mathworks_hil_bms_video}, a rate driven by the combination of the fastest ECM time constants relevant to state tracking (Chapter~2, Section~2.2.2) and the communication and control-loop timing constraints of the broader vehicle or system architecture in which the BMS operates. This 100--500~Hz range is the reference against which the per-cycle execution times reported in Section~\ref{ssec:algorithm_complexity} must be judged: an algorithm whose worst-case execution time exceeds the corresponding sample interval (2--10~ms) cannot be deployed at that rate without either algorithmic simplification or a faster target processor, regardless of its estimation accuracy in the transient-tracking evaluation of Section~\ref{sec:transient_tracking}.

\subsection{Algorithm-Level Computational Complexity}
\label{ssec:algorithm_complexity}

\subsubsection{Qualitative Complexity Comparison}

Chapters~5 through~9 developed a range of estimation algorithms -- RLS and its variants, the EKF, the UKF, particle filtering, PI and sliding-mode observers, and the multi-innovation stochastic gradient methods of Chapter~9 -- each with a distinct computational profile. Three structural factors, common across this range of methods, principally determine per-cycle computational cost:

\begin{itemize}
	\item \textbf{State/parameter dimension $n$.} Matrix operations central to RLS and Kalman-filter-type methods (covariance propagation, matrix inversion for the Kalman gain) scale, in the worst case, cubically with the dimension of the estimated state/parameter vector, making the choice of ECM order (Chapter~2, Rint versus Thevenin versus dual-polarization) a direct computational-cost lever: reported embedded implementation results show that an EKF built on a second-order RC model executes measurably slower than the same EKF architecture built on a simpler Rint model on identical target hardware \cite{energy_efficient_soc_embedded_2024}.
	\item \textbf{Sampling strategy.} The UKF's deterministic sigma-point sampling requires $2n+1$ sigma points to be individually propagated through the (generally nonlinear) process and measurement models at every time step, in contrast to the EKF's single-point linearization, which is the principal reason the UKF is consistently reported as more computationally demanding than the EKF for the same underlying state dimension \cite{kf_comparison_complexity_mdpi_2023}.
	\item \textbf{Iterative or population-based structure.} Particle filtering (Chapter~5, Section~5.3) and the heuristic optimization methods of Chapter~4 (GA, PSO) both scale with the number of particles or population members maintained, generally making them the most computationally demanding methods considered in this report when deployed with a particle/population count sufficient for good accuracy, which is part of why these methods are more commonly associated, in the literature synthesized in Chapters~4 and~5, with offline or non-time-critical online use rather than with the tightest real-time BMS execution loops.
\end{itemize}

\subsubsection{Formal Per-Cycle Operation Count}

The qualitative complexity ordering of Section~\ref{ssec:algorithm_complexity} can be stated more precisely in terms of the dominant arithmetic operation count per execution cycle, as a function of the state/parameter dimension $n$ (Table~\ref{tab:ch11_notation}). For a Kalman-filter-type method (RLS, EKF, or UKF), the covariance update step involves matrix-matrix multiplications and, for the Kalman gain calculation, a matrix inversion over the measurement dimension; where the measurement dimension is small (commonly one, for a scalar terminal-voltage measurement), this inversion reduces to a scalar division, and the dominant cost is instead the $n \times n$ covariance propagation, which scales as

\begin{equation}
	C_{\text{KF}} = \mathcal{O}(n^3),
	\label{eq:kf_complexity}
\end{equation}

\noindent reflecting the standard cubic scaling of general dense matrix-matrix multiplication with matrix dimension. For the EKF specifically, this cost is incurred once per cycle, following a single Jacobian linearization of the (generally nonlinear) process and measurement models around the current state estimate. For the UKF, the same order-$n^3$ covariance-propagation cost recurs, but it must additionally be evaluated separately for each of the $2n+1$ sigma points required by the unscented transform, giving an overall per-cycle cost of approximately

\begin{equation}
	C_{\text{UKF}} \approx (2n+1)\, \mathcal{O}(n^3) = \mathcal{O}(n^4),
	\label{eq:ukf_complexity}
\end{equation}

\noindent which formalizes, in terms of asymptotic scaling with state dimension, the empirically observed UKF-versus-EKF cost gap reported in Section~\ref{ssec:algorithm_complexity}'s embedded timing results. For a basic RLS estimator (Chapter~5, Section~5.1) operating on a parameter vector of dimension $p$ (typically smaller than the full state/parameter dimension $n$ used by a joint co-estimation architecture, Chapter~7), the corresponding per-cycle cost is $\mathcal{O}(p^2)$, reflecting the outer-product-based covariance update characteristic of the recursive least-squares update law, without the additional Jacobian-evaluation overhead of the EKF. These asymptotic relations are not, on their own, sufficient to predict absolute execution time on a specific target microcontroller -- constant factors, memory-access patterns, and compiler optimization (Section~\ref{ssec:embedded_optimization}) all materially affect the realized timing, as the reported figures in Section~\ref{ssec:algorithm_complexity} illustrate -- but they correctly predict the \emph{relative ordering} of cost across methods and across ECM order (via $n$), consistent with every reported embedded comparison cited in this section.

\begin{table}[htbp]
	\centering
	\caption{Formal per-cycle computational complexity of estimation methods, as a function of state/parameter dimension.}
	\label{tab:formal_complexity}
	\begin{tabular}{p{2.6cm}p{3.4cm}p{3.4cm}p{3.4cm}}
		\toprule
		Method & Dominant per-cycle cost & Scaling & Governing equation \\
		\midrule
		RLS & Covariance/gain update & $\mathcal{O}(p^2)$ & Sec.~5.1 update law \\
		EKF & Single-point covariance propagation & $\mathcal{O}(n^3)$ & Eq.~\eqref{eq:kf_complexity} \\
		UKF & $(2n{+}1)$-point covariance propagation & $\mathcal{O}(n^4)$ & Eq.~\eqref{eq:ukf_complexity} \\
		\bottomrule
	\end{tabular}
\end{table}

\subsubsection{Reported Embedded Timing Results}
\label{ssec:reported_embedded_timing}

Direct, hardware-measured timing comparisons corroborate the qualitative ordering above. A comparative real-time SOC estimation study using an embedded Raspberry-Pi-based test bench reported that UKF execution time per unit cell increased by approximately 7~seconds, and load rate by approximately 12.4\%, relative to an EKF implementation under the same experimental conditions, attributing the difference directly to the UKF's higher-dimensional sigma-point sampling requirement \cite{realtime_soc_embedded_electronics_2022}. A separate study targeting an STM32H7-based microcontroller reported that an EKF built on a second-order RC ECM required approximately 30\% more execution time than the same EKF architecture built on a simple Rint model \cite{energy_efficient_soc_embedded_2024}, directly illustrating the state-dimension scaling effect noted above in a genuinely embedded (rather than desktop or single-board-computer) context. A third study, focused specifically on EKF implementation optimization for an ARM Cortex-M3 target, reported reducing single-run EKF execution time from 260.4~$\mu$s to 37.7~$\mu$s -- an approximately sevenfold reduction -- through a combination of loop unrolling, function inlining, lookup-table substitution for transcendental function evaluation, and restructured data access, without measurable loss of estimation accuracy \cite{ekf_embedded_optimization_researchgate}.

\begin{table}[htbp]
	\centering
	\caption{Qualitative and reported comparative computational cost of estimation methods developed in this report.}
	\label{tab:computational_cost}
	\begin{tabular}{p{2.6cm}p{2.0cm}p{5.0cm}p{2.6cm}}
		\toprule
		Method & Relative cost & Dominant cost driver & Report reference \\
		\midrule
		RLS (Ch.~5) & Low & Single covariance/gain update per cycle, dimension-dependent & Sec.~5.1 \\
		EKF (Ch.~5) & Low--moderate & Single-point Jacobian linearization; scales with ECM order & Sec.~5.2.1 \\
		UKF (Ch.~5) & Moderate--high & $2n+1$ sigma-point propagation per cycle & Sec.~5.2.2 \\
		Particle filter (Ch.~5) & High & Scales with particle count & Sec.~5.3 \\
		PI / SMO observers (Ch.~6) & Low & Simple gain-based correction; no covariance propagation & Sec.~6.2--6.3 \\
		GA / PSO (Ch.~4) & High (offline) & Scales with population size and generation count & Sec.~4.3 \\
		\bottomrule
	\end{tabular}
\end{table}

\subsection{Embedded Optimization Techniques}
\label{ssec:embedded_optimization}

Beyond algorithm selection, several implementation-level optimization techniques, documented in the embedded-BMS literature, can substantially reduce the execution time of a given algorithm without altering its mathematical formulation, and are therefore directly applicable regardless of which estimation method from Chapters~5 through~9 is ultimately selected. It is worth noting explicitly, before cataloguing these techniques, that they operate at a level below the algorithmic complexity analysis of Section~\ref{ssec:algorithm_complexity}: none of them changes the asymptotic scaling relations of Equations~\eqref{eq:kf_complexity}--\eqref{eq:ukf_complexity}, but each reduces the constant factor multiplying that scaling, which is frequently the dominant lever available once an algorithm family has already been selected on accuracy/cost grounds (Section~\ref{ssec:realtime_tracking_tradeoff}).

\begin{itemize}
	\item \textbf{Lookup-table substitution.} Transcendental function evaluations -- most notably the $V_{oc}(z)$ curve fitted in Chapter~3, Section~3.2.3, and any exponential terms arising from the discretized RC dynamics of Chapter~2, Equation~2.9 -- are commonly replaced with pre-computed lookup tables and interpolation, avoiding repeated, comparatively expensive floating-point transcendental function calls at every execution cycle \cite{ekf_embedded_optimization_researchgate}.
	\item \textbf{Fixed-point arithmetic.} Where the target microcontroller lacks a hardware floating-point unit, or where further throughput margin is required beyond what floating-point execution provides, fixed-point implementations of the estimation algorithm's arithmetic can reduce per-operation cost substantially, at the cost of additional implementation complexity in managing numerical scaling and overflow.
	\item \textbf{Matrix operation restructuring.} For Kalman-filter-type methods, separating the computation of the Kalman gain and covariance update from the state prediction and measurement update steps, and exploiting the sparsity or structure of the specific ECM's system matrices (for instance, the diagonal structure of $A_d$ in Chapter~2, Equation~2.9), reduces the number of arithmetic operations relative to a generic, dense-matrix implementation of the same filter equations \cite{ekf_embedded_optimization_researchgate}.
	\item \textbf{Reduced-order model selection.} As already noted in Section~\ref{ssec:algorithm_complexity}, selecting the lowest ECM order (Chapter~2) that meets the transient-tracking accuracy requirements established in Section~\ref{sec:transient_tracking} is itself an effective, algorithm-independent optimization, since state/parameter dimension is a first-order driver of computational cost across every method family considered in this report.
\end{itemize}

\subsection{Automatic Code Generation and the SIL/PIL Toolchain}
\label{ssec:code_generation}

The optimization techniques of Section~\ref{ssec:embedded_optimization} are, in current model-based development practice, frequently applied not by hand-writing target C code, but through automatic code generation from the same Simulink model used for the MIL and SIL validation stages of Section~\ref{ssec:xil_taxonomy} \cite{mathworks_bms_solutions}. This toolchain continuity -- generating deployment code directly from a model that has already been validated at the MIL and SIL levels -- is part of what gives SIL testing its practical value beyond pure algorithm verification: because the SIL stage exercises the \emph{same generated code} that will ultimately run on target hardware (Section~\ref{ssec:xil_taxonomy}), the embedded optimization techniques of Section~\ref{ssec:embedded_optimization} can, in many toolchains, be configured as code-generation settings (fixed-point data typing, lookup-table approximation of transcendental functions) rather than as manual re-implementation, reducing the risk that a hand-optimized deployment implementation silently diverges, numerically, from the algorithm validated in Sections~\ref{sec:sil_coupling} and~\ref{sec:transient_tracking}. Verifying that this divergence has not occurred is, in turn, one of the specific purposes of the PIL validation stage that follows SIL in the hierarchy of Section~\ref{ssec:xil_taxonomy}: a PIL run compares the target-compiled code's output against the same SIL-level reference trajectories used throughout this chapter, isolating any accuracy loss introduced specifically by the code-generation and target-numerics step.

\subsection{Memory Footprint and Numerical Precision}
\label{ssec:memory_footprint}

Execution time is not the only computational resource constrained on an embedded BMS microcontroller; RAM and non-volatile memory footprint, and the numerical precision available for a given implementation, are equally binding constraints in practice. Kalman-filter-type methods (Chapter~5) require persistent storage of the state covariance matrix, whose memory footprint scales quadratically with state/parameter dimension $n$ -- an $n \times n$ matrix of double-precision floating-point values -- so that the same dimension-scaling argument raised for execution time in Section~\ref{ssec:algorithm_complexity} applies equally, and often more restrictively, to memory footprint on the smaller microcontroller variants common in cost-sensitive BMS designs. Observer-based methods (Chapter~6), by contrast, generally require only the observer state itself and a small number of gain constants, giving them a materially smaller memory footprint than an equivalently sized Kalman-filter implementation -- one of the practical advantages noted qualitatively for observer-based methods in Chapter~6, and now given a concrete computational-resource grounding here.

Numerical precision interacts with these memory constraints in a way that is specific to the ECM structures of Chapter~2: the RC-branch discretization of Chapter~2, Equation~2.9 involves an exponential term $e^{-\Delta t/\tau_i}$ that, for a slow time constant $\tau_i$ and a short sample interval $\Delta t$, evaluates very close to unity, so that the corresponding state-transition coefficient can lose significant precision if computed and stored in a low-precision (e.g. single-precision or fixed-point) format without care. This is a numerical, rather than algorithmic, consideration, but it is one that the embedded-optimization techniques of Section~\ref{ssec:embedded_optimization} -- particularly lookup-table substitution and fixed-point scaling -- must be designed around explicitly, since a naively chosen fixed-point scaling factor can silently degrade RC-branch tracking accuracy in a way that would not be visible in a floating-point SIL simulation of the kind developed in this chapter, but would surface only at the PIL stage (Section~\ref{ssec:xil_taxonomy}) once genuine target-hardware numerics are involved.

\subsection{Scalability to Pack-Level Estimation}
\label{ssec:pack_scalability}

Every model and estimation structure developed in this report has been scoped, following the single-cell assumption stated explicitly in Chapter~2, Section~2.1.1, to a single electrochemical cell. A production BMS, however, must run its estimation algorithm -- whichever one is selected following the accuracy/cost trade-off of Section~\ref{ssec:realtime_tracking_tradeoff} -- independently (or with limited cross-cell coupling) across every cell in a series/parallel pack, which can range from a handful of cells in a consumer device to several hundred in an electric vehicle traction pack. The per-cell computational cost figures of Section~\ref{ssec:algorithm_complexity} must therefore be scaled by the number of independently estimated cells (or cell groups, where cells are grouped and estimated collectively to reduce cost) to obtain the true aggregate computational load on the BMS's central or distributed processing hardware, a scaling relationship reported explicitly in HIL-oriented BMS validation contexts, where achieving a target real-time step becomes progressively more difficult as the number of series-connected cells being simulated or estimated grows \cite{mathworks_hil_bms_video}. This pack-level scaling consideration is noted here as a boundary of this report's single-cell scope, consistent with the scope statement of Chapter~2, Section~2.1.1, rather than developed further, since pack-level architecture (centralized versus distributed BMS topologies, cell-to-cell communication overhead) lies outside the parameter-estimation focus established in Chapter~1.

\subsection{SIL Test Harness: Algorithmic Summary}
\label{ssec:sil_harness_pseudocode}

The SIL evaluation procedure developed across Sections~\ref{sec:sil_coupling} through~\ref{sec:realtime_efficiency} can be summarized in the following pseudocode, consolidating the signal flow of Section~\ref{ssec:sil_architecture}, the test-case structure of Section~\ref{ssec:sil_test_design}, the metrics of Section~\ref{ssec:transient_metrics}, and the repeated-trial evaluation practice of Section~\ref{ssec:sil_reproducibility} into a single procedure.

\begin{quote}
\textbf{Algorithm: SIL evaluation of an estimation algorithm against the Chapter~10 plant model} \\[2pt]
\textbf{Input:} plant model (Chapter~10) with parameter set from Chapters~4--9; estimation algorithm under test; test-scenario current profiles (Section~\ref{ssec:sil_test_design}); noise model (Chapter~3, Section~3.5.2); number of trials $N_{\text{trial}}$. \\
\textbf{Output:} distribution of transient-tracking metrics (Table~\ref{tab:sil_metric_summary_template}) and per-cycle execution-time statistics.
\begin{enumerate}
	\item For each test scenario (idle/partial-load, load disconnect, high-current/thermal, drive-cycle):
	\begin{enumerate}
		\item For trial $j = 1, \dots, N_{\text{trial}}$:
		\begin{enumerate}
			\item Draw an independent noise realization from the Chapter~3 measurement-noise model.
			\item Simulate the plant model under the scenario's current profile, logging true states/parameters $z(t), U_p(t), R_0, R_1, C_1$.
			\item Apply the noise realization to the plant's terminal voltage (and current, if a current-sensor model is included) to form the estimator's measured input.
			\item Execute the compiled estimation algorithm against the noisy measurements, logging estimated states/parameters $\hat{z}(t)$, $\hat{R}_0(t)$, etc., and per-cycle execution time (host-PC SIL timing; indicative only, per Section~\ref{ssec:sil_nonrealtime}).
			\item Compute the trial's RMSE, convergence time, and overshoot metrics (Section~\ref{ssec:transient_metrics}) from the logged true and estimated trajectories.
		\end{enumerate}
		\item Aggregate the $N_{\text{trial}}$ trial-level metrics into a mean and standard deviation for the scenario (Table~\ref{tab:sil_metric_summary_template}).
	\end{enumerate}
	\item Compare the aggregated metrics across candidate estimation algorithms (e.g. EKF versus UKF versus a Chapter~6 observer) to inform the accuracy/cost selection discussed in Section~\ref{ssec:realtime_tracking_tradeoff}.
	\item Where PIL or HIL facilities are available, repeat a representative subset of scenarios at those validation levels (Section~\ref{ssec:xil_taxonomy}) to confirm that SIL-level conclusions carry over to target-representative numerics and genuine real-time execution.
\end{enumerate}
\end{quote}

\subsection{Interpreting Real-Time Results Alongside Transient Tracking}
\label{ssec:realtime_tracking_tradeoff}

The results of Sections~\ref{ssec:algorithm_complexity} and~\ref{ssec:embedded_optimization}, considered together with the transient tracking results of Section~\ref{sec:transient_tracking}, complete the accuracy-versus-cost trade-off that has recurred throughout this report since it was first raised, in general terms, in Chapter~2, Section~2.1.5 (model selection considerations) and revisited for individual algorithm families in Chapters~5 through~7: a UKF-based estimator may track the abrupt load-disconnect transient of Section~\ref{ssec:load_disconnect} with lower RMSE than an EKF-based estimator built on the same ECM, but the comparative embedded timing results of Section~\ref{ssec:algorithm_complexity} indicate that this accuracy advantage comes at a measurable, and on constrained hardware potentially prohibitive, computational cost. The SIL environment developed in this chapter (Section~\ref{sec:sil_coupling}) is, for this reason, best used not to identify a single universally preferred estimation method, but to generate the paired accuracy/cost evidence -- transient tracking metrics from Section~\ref{sec:transient_tracking} alongside the timing considerations of this section -- needed to make that selection deliberately, for a specific target microcontroller and a specific application's accuracy requirements.

\section{Summary}
\label{sec:summary_ch11}

This chapter coupled the physical plant model of Chapter~10 to the algorithmic estimation methods of Chapters~5 through~9 within a software-in-the-loop simulation environment, positioning this validation activity within the broader MIL/SIL/PIL/HIL verification hierarchy used in automotive embedded-systems practice (Table~\ref{tab:xil_levels}) and establishing the closed-loop signal architecture, including the ground-truth availability that distinguishes SIL from later, more hardware-realistic validation stages, used throughout the remainder of the chapter. Transient tracking performance (Section~\ref{sec:transient_tracking}) was evaluated using the abrupt load-disconnect scenario as the primary stress test, with the plant's expected step-plus-relaxation response derived directly from the Thevenin model of Chapter~2 and three distinct estimator-side failure modes -- delayed convergence, transient parameter corruption, and model-mismatch-induced divergence -- distinguished using the ground-truth signals unique to the SIL environment (Table~\ref{tab:load_disconnect_modes}). The related concern of spurious voltage-estimate oscillation was addressed separately, tracing the phenomenon to Kalman-gain and observer-gain tuning choices already developed in Chapters~5 and~6, and identifying three general stabilization strategies drawn from that same material. Real-time computational efficiency (Section~\ref{sec:realtime_efficiency}) then quantified the cost side of the accuracy/cost trade-off implicit throughout this report, using reported embedded-hardware timing results to corroborate the qualitative complexity ordering of RLS, EKF, UKF, particle filtering, and observer-based methods, and cataloguing implementation-level optimization techniques applicable regardless of which specific algorithm is ultimately selected.

Table~\ref{tab:ch11_consolidated} consolidates, for ease of reference, which section of this chapter addresses which combination of estimation-method family and evaluation criterion, drawing together the accuracy-side results of Section~\ref{sec:transient_tracking} and the cost-side results of Section~\ref{sec:realtime_efficiency}.

\begin{table}[htbp]
	\centering
	\caption{Consolidated index of Chapter~11 evaluation results by estimation-method family.}
	\label{tab:ch11_consolidated}
	\begin{tabular}{p{2.8cm}p{3.4cm}p{3.4cm}p{3.0cm}}
		\toprule
		Method family & Transient tracking evaluation & Real-time cost evaluation & Report chapter \\
		\midrule
		RLS / variants & Sec.~\ref{sec:transient_tracking}, general metrics & Low cost, Sec.~\ref{ssec:algorithm_complexity} & Ch.~5 \\
		EKF & Secs.~\ref{ssec:load_disconnect}--\ref{ssec:voltage_peak_stabilization} & Low--moderate, Sec.~\ref{ssec:algorithm_complexity} & Ch.~5 \\
		UKF & Secs.~\ref{ssec:load_disconnect}--\ref{ssec:voltage_peak_stabilization} & Moderate--high, Sec.~\ref{ssec:algorithm_complexity} & Ch.~5 \\
		Particle filter & Sec.~\ref{sec:transient_tracking}, general metrics & High, Sec.~\ref{ssec:algorithm_complexity} & Ch.~5 \\
		PI / SMO observers & Sec.~\ref{ssec:voltage_peak_stabilization} & Low, Sec.~\ref{ssec:algorithm_complexity} & Ch.~6 \\
		Co-estimation architectures & Sec.~\ref{ssec:load_disconnect}, mode 2 & Scales with $n$, Sec.~\ref{ssec:algorithm_complexity} & Ch.~7 \\
		\bottomrule
	\end{tabular}
\end{table}

Chapter~12 concludes this report by synthesizing the methodologies developed across all preceding chapters, revisiting the accuracy/cost/interpretability trade-offs first framed in Chapter~2 and now grounded in the transient-tracking and real-time-efficiency evidence generated in this chapter, and identifying the remaining open challenges -- in parameter identifiability, aging-aware co-estimation, and embedded deployment -- that define the road ahead for lithium-ion battery parameter estimation.
